\documentclass[leqno]{jsarticle}
\usepackage{amssymb}
\usepackage{amsthm}
\usepackage{amsmath}
\usepackage{comment}
	%可換図式用
		%\usepackage[all]{xy}
	%重くなるので使わいときはコメント化しておく。
%定理環境 thmはあまり使わずpropをなるべく使う。
%主張は基本的に証明の中で使い、そのため番号を付けない。
%注意環境を付けてみた。
\newtheorem{thm}{定理}
\newtheorem{prop}[thm]{命題}
\newtheorem{cor}[thm]{系}
\newtheorem{lem}[thm]{補題}
\newtheorem{df}[thm]{定義}
\newtheorem{exam}[thm]{例}
\newtheorem{fact}[thm]{事実}
\newtheorem*{claim}{主張}
\newtheorem*{cau}{注意}
\renewcommand{\proofname}{{\bf 証明}}
%矢印たち。
%\toは\rightarrowと同じ。\getsは\leftarrowと同じ。同値は\iff
%上向き矢はuparrow逆はdownarrow
\newcommand{\To}{\Longrightarrow}
\newcommand{\Gets}{\Longleftarrow}
%黒板太字
\newcommand{\nn}{\mathbb{N}}
\newcommand{\zz}{\mathbb{Z}}
\newcommand{\qq}{\mathbb{Q}}
\newcommand{\rr}{\mathbb{R}}
\newcommand{\cc}{\mathbb{C}}
%無限大
\newcommand{\mgn}{\infty}
%ギリシャン
\newcommand{\dpst}{\displaystyle}
\newcommand{\ep}{\varepsilon}
\newcommand{\al}{\alpha}
\newcommand{\be}{\beta}
\newcommand{\la}{\lambda}
\newcommand{\La}{\Lambda}
\newcommand{\ga}{\gamma}
%商集合
\newcommand{\qt}[2]{#1/\! #2}
%enumerate環境
\renewcommand{\labelenumi}{{\rm (\arabic{enumi})}}
%以降alignじゃなくてenumerate を使え。
%なんか演算
\newcommand{\card}{\text{  {\rm card}} }
\newcommand{\supp}{\text{  {\rm supp}} }
%なんか演算２
\newcommand{\bd}{\text{  {\rm Bdry}} }
\newcommand{\cl}{\text{  {\rm CL}} }
\newcommand{\nai}{\text{  {\rm Int}} }
\newcommand{\uu}{\mathcal{U}}
\newcommand{\vv}{\mathcal{V}}
\newcommand{\A}{\mathcal{A}}
\newcommand{\B}{\mathcal{B}}
\newcommand{\F}{\mathcal{F}}
\newcommand{\s}{\text{  {\rm St}} }
\newcommand{\C}{\mathcal{C}}
\newcommand{\D}{\mathcal{D}}
\newcommand{\E}{\mathcal{E}}
\newcommand{\W}{\mathcal{W}}
\newcommand{\dd}{\Delta}
\newcommand{\pr}{\prec}
\newcommand{\zp}{\zz_{+}}
\newcommand{\zr}{\zz_{\ge0}}
\newcommand{\wsp}{\text{{\rm wspp}}}
\newcommand{\wspp}{\text{{\rm wspp}}}
\newcommand{\lil}{\la\in \La}
\newcommand{\de}{\delta}

%差集合は\setminusで統一する。等号の否定は\neq
%真部分集合は\subsetneqq　偏微分は\partial
%ディスプレイスタイル。\dfracでディスプレイ分数
%空集合は\Oを使う！！！！！！！！！！！！

\title{パラコンパクト性：Route X}
\author{一色三良(concious77)}
\date{\today}
			\begin{document}
\maketitle
パラコンパクト性はいろいろな特徴付けが知られているが、この文書の目的はそれらの特徴付けの正規被覆への相対化である定理\ref{main}の証明である。パラコンパクトハウスドルフ空間はその任意の開被覆が正規被覆になるような空間であることに注意せよ。

\begin{thm}[Stone-Michael-Morita]\label{main}
位相空間$X$の開被覆$\uu$について以下は同値
\begin{enumerate}
\item $\uu$は正規被覆
\item $X$から距離空間$Y$への写像$f$と$Y$の開被覆$\vv$が存在し、$f^{-1}(\vv)\prec \uu$
\item $\uu$は局所有限な正規被覆開細分を持つ
\item $\uu$は局所有限かつ$\sigma$疎な正規被覆開細分を持つ
\item 連続写像の族$\{f_{\la}:X\to[0,1]\}_{\la\in \La}$が存在し$H_{\la}=\{x\ | f_{\la}(x)>0\}$とし$\mathcal{H}=\{H_{\la}\}_{\la\in \La}$としたとき、$\mathcal{H}$は局所有限で$\mathcal{H\prec \uu}$。つまり局所有限なコゼロ被覆で細分されるということ。
\item $\uu$に弱く従属する局所有限な単位の分割が存在する。
\item $\uu$に強く従属する局所有限な単位の分割が存在する。
\item $\uu$に弱く従属する単位の分割が存在する。
\item $\uu$に強く従属する単位の分割が存在する。
\item $X$の正規被覆$\A$が存在して各$S\in \A$について$\{U\cap S\ |\ U\in \uu\}$は部分空間$S$の正規被覆となる。
\end{enumerate}

\end{thm}

\section{正規被覆}

\begin{df}[星型集合]
集合$X$とその部分集合族$\mathcal{A}$、及び$S\subset X$について
\[
\s(S,\mathcal{A})=\bigcup\{A\ |\ A\in \mathcal{A},\ S\cap A\neq\O\}
\]
を$\mathcal{A}$による$S$を中心とする星型集合という。
\end{df}

\begin{df}[$*$と$\Delta$]
$X$の被覆$\uu$について$\uu^{*}$を
\[
\uu^{*}=\{\s(S,\uu)\ |\ S\in \uu\}
\]
と定義し、そして
$\uu^{\Delta}$を
\[
\uu^{\Delta}=\{\s(x,\uu)\ |\ x\in X\}
\]
と定義する。
\end{df}


\begin{df}[細分]
$\mathcal{A},\mathcal{B}を$集合$X$の部分集合の族とする。このとき$\mathcal{A}$が$\mathcal{B}$の細分であるとは
\[
\forall B\in \mathcal{B}\ \exists A\in \mathcal{A}\ A\subset B
\]
となることである。\par
$\A$が$\B$の細分であるとき
\[
\A\pr \B
\]
と書くことにしよう。
\end{df}


\begin{df}[星型細分、$\Delta$細分]
集合$X$の部分集合族$\mathcal{A},\B$について$\A$が$\B$の星型細分であるとは
\[
\A^{*}\pr \B
\]
となることである。\par
また、集合$X$の部分集合族$\mathcal{A},\B$について$\mathcal{A}$が$\B$の$\Delta$細分、又は重心細分であるとは
\[
\A^{\dd}\pr\B
\]
\end{df}

\begin{lem}\label{benri}
$A\subset B,\uu\pr \mathcal{V}$ならば
\[
\s(A,\uu)\subset \s(B,\mathcal{V})
\]
\end{lem}
\begin{proof}
まず$\s(A,\uu)\subset \s(B,\uu)$は明らかである。次に$\s(B,\uu)\subset \s(B,\vv)$を示そうとするのだが、細分なので明らかである。よって証明は終わる。
\end{proof}


\begin{lem}\label{aaa}$X$の被覆$\A,\B$について
\[
\A^{\dd}\pr \B\To \A\pr\B
\]
\end{lem}
\begin{proof}
自明
\end{proof}




\begin{lem}\label{aaa}$X$の被覆$\A,\B$について
\[
\A^{*}\pr \B\To \A^{\dd}\pr\B
\]
\end{lem}
\begin{proof}
補題\ref{benri}から分かる。
\end{proof}

\begin{lem}\label{bebe}$X$の被覆$\A,\B,\C$
\[
\A^{\dd}\pr\B,\B^{\dd}\pr\C\To \A^{*}\pr\C
\]
\end{lem}
\begin{proof}
$A\in \A$を任意に与える。$a\in A$を適当に取る。$V\in \A,A\cap V\neq\O$とすると$\A^{\dd}\pr\B$よりある$B\in\B$が存在し、$A\cup V\subset B$となる。このとき$a\in B$でもある。よってこれから
\[
\s(A,\A)\subset \s(a,\B)
\]
そして$\B^{\dd}\pr \C$なので$\A^{*}\pr \C$となることはすぐ分かる。
\end{proof}

\begin{lem}
$X$の部分集合$S$と$X$の被覆$\A,\B$について
$\A^{*}\pr \B$のとき
\[
\s(\s(S,\A),\A)\subset  \s(S,\B)
\]
が成り立つ。
\end{lem}
\begin{proof}
明らか
\end{proof}


\begin{df}[正規列、正規被覆]
$X$の開被覆の列$\{\uu_n\}_{n\in \zr}$が正規列であるとは任意の$n$について
\[
\uu_{n+1}^{\dd}\pr \uu_n
\]
となること。\par
また、開被覆$\mathcal{V}$が正規被覆であるとは正規列$\{\uu_n\}_{n\in \zr}$が存在して
\[
\uu_0\pr\mathcal{V}
\]
となること。
\end{df}


\setcounter{equation}{0}
\begin{lem}\label{ccc}
$X$の被覆$\mathcal{V}$が正規被覆であることと、被覆の列$\{\uu_n\}_{n\in \zr}$が存在して
\[
\uu^{*}_{n+1}\pr\uu_n,\ \uu_0\pr\mathcal{V}
\]
を充たすことは同値である。

\end{lem}
\begin{proof}
補題\ref{aaa},\ref{bebe}からすぐに分かる。
\end{proof}


\begin{lem}\label{hey}
位相空間$X$の開被覆$\uu$と集合$A$と開集合$U$について
\[
\s(A,\uu)\subset U
\]
が成り立つならば$\cl(A)\subset U$となる。
\end{lem}
\begin{proof}
$a\in \cl(A)$を任意に与えると$\uu$が被覆であることから$V\in \A,a\in V$となる$V$が存在するが、$a$は$A$の触点なので$V\cap A\neq \O$である。よって$V\subset \s(A,\uu)\subset U$である。故に$a\in U$である。以上から$\cl(A)\subset U$が分かる。
\end{proof}

\setcounter{equation}{0}
\begin{thm}[Stone]\label{stone}
位相空間$X$の正規被覆には局所有限な開細分被覆が存在する。
\end{thm}
\begin{proof}
%%%%%%%%%%%%%
$\uu$を正規被覆とした時、補題\ref{ccc}より星型細分の列が存在するが、適当に番号を付け替えることによって開被覆$\vv$と開被覆列$\{\A_n\}_{n\in \zr}$が存在して
\[
\A_0\pr \vv,\ \vv^{*}\pr\uu
\]
でかつ
\[
\A_{n+1}^{*}\pr\A_n
\]
を充たす。\par
そして帰納的に
\[
\B_1=\A_1,\ \B_2=\{\s(V,\A_1)\ |\ V\in \B_1\},\dots,\B_n=\{\s(V,\A_n)\ |\ V\in \B_{n-1}\},\dots
\]
と置く。この時各$n\in \zp$について$\{\s(V,\A_n)\ |\ V\in\B_n\}$は$\A_0$の細分である。このことを帰納法によって証明しよう。まず$n=1$の時は明らかである。次に$n=k$までは成り立っているとして$n=k+1$のとき、$V\in \B_{k+1}$を与え$U=\s(V,\A_{k+1})$を考える。この時$V$は$\B_{k+1}$の作り方から$S\in \B_k$を用いて$V=\s(S,\A_{k+1})$と表される。そして$\A_{k+1}^{*}\pr \A_k$なので
\[
U=\s(V,\A_{k+1})=\s(\s(S,\A_{k+1}),\A_{k+1})\subset \s(S,\A_k)
\]
なので帰納法の仮定から$k+1$の場合も成り立つ。\par
以上のことから特に各$\B_n$は$\A_0$の細分である。\par
%%第二弾楽
次に$X$を整列させてその順序を$<$とする。そして$(n,x)\in \zp\times X$について
\[
C_n(x)=\s(x,\B_n)\setminus \bigcup_{y<x}\s(y,\B_{n+1})
\]
と定義する。この時$\C=\{C_n(x)\ |\ n\in \zp,x\in X\}$は$X$の被覆となる。実際$p\in X$について
\[
P=\{z\in X\ |\ p\in \bigcup_{n\in \zp}\s(z,\B_n)\}
\]
とを考えると$p\in P$なので$P$は空では無いので最小限$y$が存在する。$y$についてある$n$が存在して$p\in \s(y,\B_n)$となる。このような$n$について$p\in C_n(y)$となる。よって$\C$は被覆となる。\par
次の主張が成り立つ。
%%%%%%%%%%%%%%%%%%%%%%%%%%
\begin{claim}[a]
$n\in \zp$を固定したとき、$U\in \A_{n+1}$は$\{C_n(x)\ |\ x\in X\}$の高々１個の元としか交わらない
\end{claim}
{\bf [主張(a)の証明]}\par
いま、$U\cap C_n(x)\neq\O$としよう。すると$x\in V$となる$V\in \B_n$が存在して$U\cap V\neq\O$よって
\[
x\in U\cup V\subset \s(V,\A_{n+1})\in \B_{n+1}
\]
故に$U\subset \s(x,\B_{n+1})$よって$\C_n(x)$の作り方から$x$を$U\cap C_n(x)\neq\O$となる最小の元とすれば$U$は$C_n(x)$以外の$\{C_n(x)\ |\ x\in X\}$の元と交わらない。\par
{\bf [主張(a)の証明終わり]}\par
%%%%%%%%%%%%%%%%%%%%%%%%%
また、$C_n(x)\subset \s(x,\B_{n})$であり、$\B_n\pr \A_0$でかつ$\A_0^{*}\pr \vv$なので$\C\pr\vv$である。\footnote{証明が終わったあと詳しく説明する}\par
%%%%%%%%%%%%%%%%%%%%%%%%%%
さて$(n,x)\in \zp\times X$について
\[
D_n(x)=\s(C_n(x),\A_{n+3}),\ E_n(x)=\s(C_n(x),\A_{n+2})
\]
と定義し、
\[
\D=\{D_n(x)\ |\ (n,x)\in \zp\times X\},\ \E=\{E_n(x)\ |\ (n,x)\in \zp\times X\}
\]
と定義する。\par
このとき
\[
\s(D_n(x),\A_{u+3})=\s(\s(C_n(x),\A_{n+3}),\A_{n+3})\subset \s(C_n(x),A_{n+2})=E_n(x)
\]
が成り立つので補題\ref{hey}から$\cl(D_n(x))\subset E_n(x)$である。\par
ここで次の主張が成り立つ。
%%%%%%%%%%%%%%%%%%%%%%%%%%%%%%%%
\begin{claim}[b]
$n\in \zp$を固定すると$U\in \A_{n+2}$は$\{E_n(x)\ |\ x\in X\}$の高々１個の元とし交わらない。特に$\{E_n(x)\ |\ x\in X\}$は疎である。
\end{claim}
{\bf [主張(b)の証明]}\par
$U\cap E_n(x)\neq\O$とすると$E_n(x)$の定義から$V\cap C_n(x)\neq\O,\ U\cap V\neq\O$となる$V\in \A_{n+2}$が存在するが、$\uu_{n+2}^{*}\pr \uu_{n+1}$なので$U\cup V\subset O$となる$O\in \A_{n+1}$が存在し、主張$(a)$から$O$は$C_n(x)$以外の$\{C_n(x)\ |\ x\in X\}$の元と交わらないので、$U$は$\{E_n(x)\ |\ x\in X\}$の高々１個の元としか交わらない。\par
{\bf [主張(b)の証明終わり]}\par
%%%%%%%%%%%%%%%%%%%%%%%%%%%%%
また、$D_n(x)\subset E_n(x)$なので主張(b)から$n$を固定すると$\{D_n(x)\ |\ x\in X\}$は疎である。よって
\[
F_n=\bigcup_{x\in X}\cl(E_n(x))
\]
は閉集合である。\par
%%%%%%%%%%%%%%%%%%%%%%%%%%%%%
そして$\C\pr \vv,\ \A_n\pr \vv$で$\vv^{*}\pr\uu$なので$\mathcal{E}\pr \uu$である。\footnote{証明が終わった後詳しく説明する。}\par
%%%%%%%%%%%%%%%%%%%%%%%%%
以上を元に$(n,x)\in \zp\times X$について
\[
W_n(x)=E_n(x)\setminus \bigcup_{i<n}F_i
\]
と定義する。$\mathcal{W}=\{W_n(x)\ |\ (n,x)\in \zp\times X\}$が$\uu$の局所有限な開細分被覆であることを示そう。\par
まず$\mathcal{E}\pr \uu$であり、$W_n(x)\subset E_n(x)$なので$\mathcal{W}\pr \uu$である。\par
次に$\mathcal{W}$が被覆であることを示そう。$p\in X$を任意に与える。
\[
\exists x\in X\ p\in \cl(D_n(x))
\]
を成り立たせる最小の$n$を$m$とすると$y\in X$が存在して$p\in \cl(D_m(y))$となり、$m$の最小性から\\
$p\not\in F_1\cup F_2\cup\cdots \cup F_{m-1}$で、かつ$p\in E_m(y)$となる。よって$p\in W_m(y)$なので$\mathcal{W}$は$X$の被覆となる。\par
最後に$\mathcal{W}$の局所有限性を示そう。$p\in X$を任意に与える。すると$p\in C_i(y)$となる$C_i(y)\in \C$が存在する。このとき
\[
\s(p,\A_{i+3})\subset D_i(y)
\]
なので$\s(p,\A_{i+3})\subset F_i$よって$k>i$ならば$\s(p,\A_{i+3})$は$W_k(x)\ (x\in X)$と交わらない。\par
そして$k< i$の時には$\A_i^{*}\pr \A_k$であり、特に$\A_k^{\dd}\pr \A_i$
なので$\s(p,\A_{i+3})\subset V_{k}$となる
$V_k\in \A_{k}$が存在する。以上と主張$(b)$を踏まえると$p\in N$となる$N\in \A_{i+3}$を適当に取れば、$N\subset V_k$なので
$N$は$\mathcal{W}$の高々$i$個の元としか交わらない。\par
これで証明が終わった。
\end{proof}
\begin{cau}[証明の中で後で説明すると言った箇所]${}$\\
{\bf [$C_n(x)\subset \s(x,\B_{n})$であり、$\B_n\pr \A_0$で$\A_0^{*}\pr \vv$なので$\C\pr\vv$である。]}の説明：$x\in M$となる$M\in \A_0$を取ると補題\ref{benri}から
\[
C_n(x)\subset \s(x,\B_{n})\subset \s(M,\A_n)\in \A_0^{*}
\]
$\A_0^{*}\pr\vv$なので$\s(M,\A_n)\subset O$となる$O\in \vv$が存在するので結局
\[
C_n(x)\subset O
\]
となる$O\in \vv$が存在することになる。よって$\C\pr \vv$


{\bf [$\C\pr \vv,\ \A_n\pr \vv$で$\vv^{*}\pr\uu$なので$\mathcal{E}\pr \uu$である。]}の説明：
$E_n(x)=\s(C_n(x),\A_{n+2})$である。$\C\pr \vv$から$C_n(x)\subset V$となる$V\in \vv$が存在する。これと$\A_{n+2}\pr \vv$と補題\ref{benri}から
\[
E_n(x)=\s(C_n(x),\A_{n+2})\subset \s(V,\vv)\in \vv^{*}
\]
である。また、$\vv^{*}\pr \uu$なので$\s(V,\vv)\subset O$となる$O\in \uu$が存在するので結局
\[
E_n(x)\subset O
\]
となる$O\in \uu$が存在することになる。よって$\E\pr \uu$
\end{cau}

%%%%%%%%%%%%%%%%%%%%%%%%%%%%%%%%%%%%%%%%%%%%%%%%%%%%%%%%%%%%%%%%%%%%%
%%%%%%%%%%%%%%%%%%%%%%%%%%%%%%%%%%%%%%%%%%%%%%%%%%%%%%%%%%%%%%%%%%%
%%%%%%%%%%%%%%%%%%%%%%%%%%%%%%%%%%%%%%%%%%%%%%%%%%%%%%%%%%%%%%%%%%
\section{距離空間のパラコンパクト性}
\setcounter{equation}{0}明らかにコンパクト空間はパラコンパクトである。さらに
パラコンパクトである空間の重要な例として距離空間がパラコンパクトであることを示そう。
以下の定理の証明の中では$0\not\in \nn$とし、$B(x,r)$で擬距離空間の$x$を中心とする$r$開球を表すとする。
\begin{thm}[Stone の定理]
擬距離空間の任意の開被覆に対し局所有限でかつ$\sigma$疎であるような開細分が存在する。特に擬距離空間はパラコンパクトである。\footnote{擬距離とは非退化性$d(x,y)=0\To x=y$を仮定しない距離のことである。}
\end{thm}
\begin{proof}M.E.Rudinによる簡潔な証明\cite{5}を引用しよう。\par 
$X$を考えている擬距離空間として
$\{C_{\al}\}_{\al\in A}$を$X$の開被覆とする。ただし添字$A$は整列されているものとする。整列可能定理を用いればこのように仮定しても一般性を失わない。そして自然数$n\in \nn$に対して開集合の族$\mathcal{U}_n=\{D_{\al,n}\}_{\al\in A}$を$n$について帰納的に構成する。今、$n$より小さい自然数については構成が終わっているとして$I_{\al,n}\subset X$を次の３つの条件を充たす$X$の点$x$全体の集合として定義する。
\begin{eqnarray}
&&\al =\min \{\be \ |\ x\in C_\be \}\\
&&\forall j<n\ x\not\in \bigcup_{\be \in A}D_{\be,j}\\ 
&&B(x,3\cdot2^{-n})\subset C_{\al}
\end{eqnarray}
そして
\[
D_{\al,n}=\bigcup_{x\in I_{\al,n}}B(x,2^{-n})
\]
と定義する。こうして
\[
\mathcal{U}_n=\{D_{\al,n}\}_{\al\in A}
\]
が定義される。$n=1$のときでも$(2)$が空な条件となるだけで、$I_{\al,1}$が構成され$\mathcal{U}_1$がきちんと定義される。以上により帰納的に開集合の族の列
\[
\mathcal{U}_1,\mathcal{U}_2,\cdots ,\mathcal{U}_n,\cdots
\]
が構成される。そして
\[
\mathcal{U}=\{D_{\al,n}\}_{(\al,n)\in A\times \nn}
\]
と置く。明らかに$D_{\al,n}\subset C_{\al}$なので$\mathcal{U}$が$\{C_{\al}\}_{\al\in A}$の細分になってる事はすぐにわかる。\par 
$\mathcal{U}$が$X$の被覆となることを見よう。$x\in X$について$x\in C_{\al}$となる最小の$\al$が存在する。そして$n\in\nn$を十分大きく取れば$B(x,3\cdot s^{-n})\subset C_{\al}$となる。この$n$についてもし
\[
\exists j<n\ x\in \bigcup_{\be \in A}D_{\be,j}
\]
ならばもちろん$x\in \bigcup \mathcal{U}$である。一方
\[
\forall j<n\ x\not\in \bigcup_{\be \in A}D_{\be,j}\\ 
\]
ならば$x\in I_{\al,n}$となるので$x\in B(x,2^{-n})\subset D_{\al,n}$つまり$x\in \bigcup\mathcal{U}$
となる。いづれにせよ$x\in \bigcup\mathcal{U}$なので
\[
X=\bigcup\mathcal{U}
\]
となり、$\mathcal{U}$が$X$の開被覆であることがわかる。\par
$\mathcal{U}$が局所有限である事を見よう。$x\in X$を任意に与えると$\mathcal{U}$が被覆なので
$x\in D_{\al,n}$となる$D_{\al,n}\in \mathcal{U}$が存在し、$j\in \nn$を十分大きく取れば
\[
B(x,2^{-j})\subset D_{\al,n}
\]
となる。このとき次が成り立つ
\begin{eqnarray*}
&&(a):i\ge n+jならば任意の\be \in AについてB(x,2^{-n-j})はD_{\be,i}と交わらない。\\
&&(b):i<n+jならばD_{\be,i}がB(x,2^{-n-j})と交わるような\be\in A は高々一個しか存在しない。
\end{eqnarray*}
この$(a),(b)$から$\mathcal{U}$の局所有限性はすぐに従う。またこの$(a),(b)$から$\mathcal{U}_{\al,n}$が疎であることもわかる。
これらの証明をしよう。\par
$(a)$の証明：$i\ge n+j$ならば特に$i>n$なので$I_{\be,i}$の定義の$(2)$から$y\in I_{\be,i}$ならば$y\not\in D_{\al,n}$そして$B(x,2^{-j})\subset D_{\al,n}$から$y\not\in B(x,2^{-j})$なので
\[
d(x,y)\ge 2^{-j}
\]
が成立する。いまここで$(a)$が成立しない、つまり$B(x,2^{-n-j})$が$D_{\be,i}$と交わると仮定すると
\[
D_{i,\be }=\bigcup_{y\in I_{i,\be}}B(y,2^{-i})
\]
よりある$y\in I_{i,\be}$について$z\in B(x,2^{-n-j})\cap B(y,2^{-i})$が存在するが、$i\ge n+j\ge n+1$及び$n+j\ge j+1$から$2^{-i}\le 2^{-j-1}$と$2^{-n-j}\le 2^{-j-1}$となるので
\[
2^{-j}\le d(x,y)\le d(x,z)+d(z,y)<2^{-n-j}+2^{-i}\le 2^{-j-1}+2^{-j-1}=2^{-j}
\]
がわかり$2^{-j}<2^{-j}$となるから矛盾。よって$(a)$が成立する。[$(a)$の証明終わり]\par
$(b)$の証明：まず$\be\neq\ga,p\in D_{\be,i},q\in D_{\ga,i}$ならば$2^{-n-j+1}<d(p,q)$が成立することを示そう。$\be<\ga$と仮定しても一般性を失わない。$D$の定義から
\[
p\in B(y,2^{-i}),q\in B(z,2^{-i})
\]
となる$y\in I_{\be,i},z\in I_{\ga,i}$が存在する。さて$\be<\ga$より$I_{\ga,i}$の定義の$(1)$から$z\not\in C_{\be }$が成り立ち、$D_{\be,i}\subset C_{\be}$から
\[
z\not\in D_{\be,i}
\]
となる。
よって$I_{\be,i}$の定義の$(3)$から$B(y,3\cdot 2^{-i})\subset D_{\be i}$なので、上と合わせて$z\not\in B(y,3\cdot 2^{-i})$であるから
\[
3\cdot 2^{-i}\le d(y,z)
\]
を得る。
これと三角不等式から
\[
3\cdot 2^{-i}\le d(y,z)\le d(y,p)+d(p,q)+d(q,z)<2^{-i}+d(p,q)+2^{-i}=2\cdot 2^{-i}+d(p,q)
\]
なので結局
\[
2^{-i}<d(p,q)
\]
そして$i<n+j$から$i\le n+j-1$なので$2^{-i}\ge 2^{-n-j+1}$となり
\[
2^{-n-j+1}<d(p,q)
\]が成り立つ。
さてもし$(b)$が成り立たない、つまり$B(x,2^{-n-j})$が$D_{\be,i},D_{\ga,i}(\be\neq \ga)$の二つと交わるなら\\$p\in B(x,2^{-n-j})\cap D_{\be,i},\ q\in B(x,2^{-n-j})\cap D_{\ga,i}$となる$p,q$が存在するが
\[
d(p,q)\le d(p,x)+d(x,q)<2^{-n-j}+2^{-n-j}=2^{-n-j+1}
\]
となるが$2^{-n-j+1}<d(p,q)$なので矛盾する。よって$(b)$が成立する。[$(b)$の証明終わり]\par
以上で証明は完結した。
\end{proof}

%%%%%%%%%%%%%%%%%%%%%%%%%%%%%%%%%%%%%%%%%%%%%%%%%%%%%%%%%%%%%%%%%%%%%%%%%%%%%%%%
%%%%%%%%%%%%%%%%%%%%%%%%%%%%%%%%%%%%%%%%%%%%%%%%%%%%%%%%%%%%%%%%%%%%%%%%%%%%%
%%%%%%%%%%%%%%%%%%%%%%%%%%%%%%%%%%%%%%%%%%%%%%%%%%%%%%%%%%%%%%%%%%%%
\section{単位の分割}
\begin{df}[細分]$\mathfrak{A},\mathfrak{B}を$集合$X$の部分集合の族とする。このとき$\mathfrak{A}$が$\mathfrak{B}$の細分であるとは
\[
\forall B\in \mathfrak{B}\ \exists A\in \mathfrak{A}\ A\subset B
\]
となることである。
\end{df}

\begin{df}[点有限族]集合$X$の部分集合の族$\mathfrak{A}$が点有限であるとは$X$の任意の点$x$に対して$x\in A$となる$A\in \mathfrak{A}$が有限個しかないことである。

\end{df}

\begin{df}[局所有限族]位相空間$X$に於いて$X$の部分集合の族$\mathfrak{A}$が局所有限であるとは$X$の各点$x$について$x$の近傍$N$が存在して
\[
N\cap A\neq \emptyset となるA\in \mathfrak{A}は有限個
\]
を充たすことである。
\end{df}

\begin{df}[開収縮]位相空間$X$の開被覆$\mathcal{U}$に対して$\{V_U\ |\ U\in \mathcal{U}\}$が開収縮であるとは$\{V_U\ |\ U\in \mathcal{U}\}$が$X$の開被覆であって、任意の$U\in \mathcal{U}$について
\[
CL(V_U)\subset U
\]
を充たすこと。
\end{df}

\begin{df}[閉収縮]位相空間$X$の開被覆$\mathcal{U}$に対して$\{F_U\ |\ U\in \mathcal{U}\}$が閉収縮であるとは$\{F_U\ |\ U\in \mathcal{U}\}$が$X$の閉被覆であって、任意の$U\in \mathcal{U}$について
\[
F_{U}\subset U
\]
を充たすこと。
\end{df}


\begin{df}[Weak supportとSupport]
関数$fX\to [0,1]$の{\bf 弱台}$\wspp (f)$を
\[
\wspp (f)=\{x\in X\ |\ f(x)>0\}
\]
と定義する。\par
関数$fX\to [0,1]$の{\bf 台}$\supp (f)$を
\[
\supp (f)=CL(\supp(f))=CL(\{x\in X\ |\ f(x)>0\})
\]

\end{df}

以下様々な種類の単位の分割を定義する。
%弱い単位の分割
\begin{df}[弱い単位の分割]位相空間$X$上の実数値連続関数の族$\{f_i\}_{i\in I}$が$X$の開被覆$\mathcal{U}$に従属する弱い単位の分割であるとは$\{\wspp f_i\}_{i\in I}$が$\mathcal{U}$の細分であり
\begin{align*}
&\sum_{i\in I} f_i=1
\end{align*}
を充たすこと
\end{df}
%強い単位の分割
\begin{df}[強い単位の分割]位相空間$X$上の実数値連続関数の族$\{f_i\}_{i\in I}$が$X$の開被覆$\mathcal{U}$に従属する強い単位の分割であるとは$\{\supp f_i\}_{i\in I}$が$\mathcal{U}$の細分であり
\begin{align*}
&\sum_{i\in I} f_i=1
\end{align*}
を充たすこと
\end{df}
%点有限な弱い単位の分割
\begin{df}[点有限な弱い単位の分割]位相空間$X$上の実数値連続関数の族$\{f_i\}_{i\in I}$が$X$の開被覆$\mathcal{U}$に従属する弱い単位の分割であるとは$\{\wspp f_i\}_{i\in I}$が$\mathcal{U}$の細分であり
\begin{align*}
&\sum_{i\in I} f_i=1\\
&\text{$\{\wspp f_i\}_{i\in I}$は点有限}
\end{align*}
を充たすこと
\end{df}

%点有限な強い単位の分割
\begin{df}[点有限な弱い単位の分割]位相空間$X$上の実数値連続関数の族$\{f_i\}_{i\in I}$が$X$の開被覆$\mathcal{U}$に従属す強い単位の分割であるとは$\{\supp f_i\}_{i\in I}$が$\mathcal{U}$の細分であり
\begin{align*}
&\sum_{i\in I} f_i=1\\
&\text{$\{\supp f_i\}_{i\in I}$は点有限}
\end{align*}
を充たすこと
\end{df}

%局所有限な弱い単位の分割
\begin{df}[局所有限な弱い単位の分割]位相空間$X$上の実数値連続関数の族$\{f_i\}_{i\in I}$が$X$の開被覆$\mathcal{U}$に従属する局所有限な弱い単位の分割であるとは$\{\wspp f_i\}_{i\in I}$が$\mathcal{U}$の細分であり
\begin{align*}
&\sum_{i\in I} f_i=1\\ 
&\text{$\{\wspp f_{i}\}_{i\in I}$が局所有限}
\end{align*}
を充たすこと
\end{df}

%局所有限な強い単位の分割
\begin{df}[局所有限な強い単位の分割]位相空間$X$上の実数値連続関数の族$\{f_i\}_{i\in I}$が$X$の開被覆$\mathcal{U}$に従属する局所有限な強い単位の分割であるとは$\{\supp f_i\}_{i\in I}$が$\mathcal{U}$の細分であり
\begin{align*}
&\sum_{i\in I} f_i=1\\ 
&\text{$\{\supp f_{i}\}_{i\in I}$が局所有限}
\end{align*}
を充たすこと
\end{df}

\begin{comment}
\begin{df}[広義の単位の分割]
関数族$\{f_{\la}\}_{\lil}$が$X$の開被覆$\mathcal{U}$に従属する単位の分割であるとは$\{\wsp(f_{\la})\}$が$\mathcal{U}$の細分になっており
\[
\sum_{\lil}f_{\la}\equiv 1
\]
を充たすときに言う。
\end{df}

\begin{df}[広義の局所有限な単位の分割]
関数族$\{f_{\la}\}_{\lil}$が$X$の開被覆$\mathcal{U}$に従属する単位の分割であるとは$\{\wsp(f_{\la})\}$が局所有限で$\mathcal{U}$の細分になっており
\[
\sum_{\lil}f_{\la}\equiv 1
\]
を充たすときに言う。
\end{df}
\end{comment}

\begin{lem}
$\{f_{\la}\}_{\lil}$が弱い単位の分割とすると、任意の$x\in X$及び任意の$\ep>0$に対して$\La$の有限集合$A$が存在して
\[
\sum_{\la\in A}f_{\la}(x)>1-\ep
\]
が成り立つ。このとき$x\in \{f_{b}>\ep\}$ならば$b\in A$が成り立つ。
\end{lem}

\begin{lem}$\{f_{\la}\}_{\lil}$が弱い単位の分割とし、
任意の$\de>0$について
\[
V_a^{\de}=\{f_a>\de\}
\]
と定義すると$\{V_{\la}^{\de}\}_{\lil}$は$X$で局所有限である。
\end{lem}
\begin{proof}
$p\in X$と$\de$について先の補題の有限集合$A$を取り
\[
N=\{\sum_{\la\in A}f_{\la}>1-\de\}
\]
を考えると$N\cap V_{a}^{\de}\neq \O$ならば先の補題から$a\in A$なので$N$は族の有限個の元としか交わらない。
\end{proof}

\begin{lem}
弱い単位の分割$\{f_{\la}\}_{\lil}$について
\[
F=\sup_{\lil}f_{\la}
\]
は連続である。
\end{lem}
\begin{proof}
$p\in X$を与える。このときある$a\in \La$について$f_a(p)>0$なので適当に$f_a(p)>\de$となる$\de>0$を取ると
$\{V_{\la}^{\de}\}_{\lil}$は先の補題から局所有限である。その局所有限性を担保する点$p$における近傍を$N$とし、有限集合$A$は
\[
N\cap V_{\la}^{\de}\neq \O\To \la \in A
\]
となるものとする。このとき$N$上で
\[
F=\max_{a\in A}f_{a}
\]
なので$F$は$N$上で連続である。よって$F$は$X$で連続である。
\end{proof}

\begin{lem}\label{unit}
位相空間$X$とその開被覆$\mathcal{U}$について以下は同値
\begin{align}
&\mbox{$\mathcal{U}$に従属する弱い単位の分割が存在する}\\
&\mbox{$\mathcal{U}$に従属する局所有限な弱い単位の分割が存在する}
\end{align}
\end{lem}
\begin{proof}
$\{f_{\la}\}_{\lil}$
先と同様に
\[
F=\sup_{\lil}f_{\la}
\]
と置く。
そして
\[
g_{\la}=\max\{f_{\la}-\frac{1}{2}F,\ 0\}
\]
と定義する。このとき$\{\wsp(g_{\la})\}_{\lil}$は$X$の被覆となる。これが局所有限であることを示そう。$p\in X$を与えると$F(p)>0$なので
\[
\frac{1}{2}F(p)>\ep
\]
となる$\ep>0$が存在する。そして有限集合$A$で
\[
\sum_{\la \in A}f_{\la}(p)>1-\ep
\]
を充たすものとする。さて
\[
N=\{\sum_{\la \in A}f_{\la}>1-\ep\}\cap \{\frac{1}{2}F>\ep\}
\]
と定義する。すると$N$は開集合で$p\in N$である。ところで
\[
N\cap \wsp(g_{\la})\neq \O
\]
とする。$x\in N\cap \wsp(g_{\la})$を取ると$g_{\la}(x)>0$なので$f_{\la}(x)>\frac{1}{2}F(x)>\ep$となる。よって$f_{\la}(x)>\ep$なので$\la \in A$となるから$N\cap \wsp(g_{\la})\neq \O$となる$\la$は有限個しかない。よって$\{\wsp(g_{\la})\}_{\lil}$は局所有限である。よって
\[
G=\sum_{\la\in \La}g_{\la}
\]
が定義でき、これは連続である。さて
\[
h_{\la}=\frac{g_{\la}}{G}
\]
と置けば
$\wsp(h_{\la})=\wsp(g_{\la})\subset \wsp(f_{\la})$なので確かに$\{h_{\la}\}_{\lil}$は$\mathcal{U}$に従属する局所有限な弱い単位の分割となる。
\end{proof}













	
\begin{thebibliography}{9}
\bibitem{1}

\end{thebibliography}




			\end{document}



\begin{comment}
[可換図式]
\[\xymatrix{
&   & (2) \ar@{=>}[rd]&  &  &  & (5) \ar@{=>}[rd]&\\ 
& (1)\ar@{=>}[ru] \ar@{=>}[rd]&  &(4)\ar@{=>}[r]&(7)& (1)\ar@{=>}[ru] \ar@{=>}[rd]& & (7)&\\ 
&   & (3) \ar@{=>}[ur]&  &  &  &(6) \ar@{=>}[ur]&
}\]

[\bigin \endを使わない方法]

{\prop[aa] aaa}
で
\begin{prop}[aa]
aaa
\end{prop}
と同じ\df \thm \proof でも同様

[直和]
$\amalg$

[同値関係でよく使う波線]
\sim

[引数付きの新命令]
\newcommand{\prn}[1]{\|#1\|_p}

[番号のリセット]

\setcounter{equation}{0}


[字体の変更]

{\rm hogehoge}と使う。

\rm ローマン
\it イタリック
\sl 斜体

[supp]

\newcommand{\supp}{\text{{\rm supp}}}

[中央揃えの改行]

	\begin{eqnarray*}
		hoge1&=&hogehoge
		hoge2&=&hogehofe

	\end{eqnarray*}

&&の部分で揃えられる。






[場合分け]

	\begin{cases}
		hoge1 & \text{状況１}\\
		hoge2 & \text{状況２}\\
		・
		・
		・
	\end{cases}



\end{comment}





